\section{Conclusion}
Limited execution memory creates a joint allocation and alphabet-design problem: participation and realization limits determine which terminal lotteries implement an accepted obligation. At saturation, the contractual tail selects an off-cap highest level, and its crossing-difference identity enables classical linear-time matrix search after continuous minimization. Away from saturation, the eligibility-prefix reduction preserves endogenous branch promises, heterogeneous intermediate costs, and actual selected-level charges in one formulation.

The price decomposition turns that geometry into an algorithmic advantage. Its ordered path recurrence solves each relaxed shared-book problem in polynomial work with variable branch count and alphabet budget, and gives exact completion bounds for prefix search. The resulting global interval is backed by a feasible lottery and a checkable exhaustive partition, not merely by source provenance or an inner multiplier. A uniform rounding theorem transfers even an interrupted catalog search to continuous bounds, and a fixed-input construction proves its mesh rate sharp.

The general concave model has exact fixed-book structure, a scalar-oracle price decomposition, and feasible discretization bounds; exact unrestricted continuous optimization is established for the rational quadratic subclass. The continuous face algorithm remains an exponential reference, and global prefix search also has exponential worst-case size. The price subproblem and verified completion bounds do not assert polynomial global optimization at the original budget. The two-price recovery theorem adds a distinct polynomial guarantee: a fixed union alphabet of at most twice the size achieves any prescribed additive accuracy when charges vanish, and an explicitly computed excess-charge penalty when they do not. Promise and realization feasibility remain exact. Feasible mesh transport extends this resource tradeoff to continuous design; it is not an FPTAS at unchanged memory. These distinctions identify the contribution's scope without replacing the original continuous-design result.

The retained phase diagram, institutional-gap bounds, repeated-cap aggregation, and robustness of off-cap design describe concrete decisions within the stipulated model. The additional monotone-charge theorem shows why a one-branch pathwise-randomization separation needs the deliberately nonmonotone charge in the retained example. The abstract ordered-interface theorem makes the method reusable independently of the operational narrative; the synthetic experiments test the mathematics and algorithms rather than claiming a calibrated deployment. Post-draw exit, hidden reports, coupled opening charges, and additional history channels remain different model specifications, not tacit consequences of this one.
